% Bagian: propositional-logic

\documentclass[../../include/open-logic-part]{subfiles}

\begin{document}

\olpart{pl}{Logika Proposisi}

\begin{editorial}
  Bagian ini memuat materi tentang logika proposisi klasik. Bab
  pertama relatif mendasar dan hanya mencantumkan definisi serta
  hasil; banyak bukti tidak dijabarkan, tetapi diserahkan sebagai
  latihan. Materi mengenai sistem bukti dan teorema kelengkapan
  diimpor dari bagian tentang logika orde pertama, dengan tag
  ``FOL'' disetel ke salah. Hal ini meniadakan segala sesuatu yang
  berkaitan dengan predikat, suku, dan kuantor, serta mengganti
  pembahasan tentang !!{structure}s~$\Struct{M}$ dengan pembahasan
  tentang !!{valuation}s~$\pAssign{v}$.

  Bagian ini direncanakan untuk diperluas agar memuat lebih banyak
  rincian serta topik dan hasil tambahan, seperti kelengkapan
  fungsional-kebenaran.
\end{editorial}

\olimport[syntax-and-semantics]{syntax-and-semantics}

\tagfalse{FOL}

\olimport[../first-order-logic/proof-systems]{proof-systems}

\iftag{prfSC}{%
  \olimport[../first-order-logic/sequent-calculus]{sequent-calculus}
}{}

\iftag{prfND}{%
  \olimport[../first-order-logic/natural-deduction]{natural-deduction}
}{}

\iftag{prfTab}{%
  \olimport[../first-order-logic/tableaux]{tableaux}
}{}

\iftag{prfAX}{%
  \olimport[../first-order-logic/axiomatic-deduction]{axiomatic-deduction}
}{}

\olimport[../first-order-logic/completeness]{completeness}

\tagtrue{FOL}

\OLEndPartHook

\end{document}
